\chapter{引言}

%标准模型，重味物理，粲，单举衰变
数千年来，人类对无穷的探索从未止步。从日心说到宇宙微波背景，从物性论到粒子物理标准模型，人类对于自然世界的认识不断深入。在这个过程中，粒子物理学作为物理学的基石之一，着眼于极小尺度的现象和其背后的规律。物质的基本组成单元是什么，这依然是现代粒子物理学需要回答的问题。

二十世纪, 科学发展的速度超过了人类历史上以往任何一个时期。在这个时期，人类对物质的认识发生了翻天覆地的变化。人类发现了电子、原子核、质子和中子等\cite{stoney1881lii, rutherford2012scattering, rutherford2010collision, chadwick1932possible}。从宇宙射线和加速器实验中发现了一系列新的粒子，比如$\mu$ 子，$\pi$ 介子，$K$ 介子等\cite{anderson1932energies,anderson1936cloud, anderson1952total, cowan1956detection, danby1962observation}。


随着对电磁相互作用\ct{schwinger1948quantum01, schwinger1949quantum, schwinger1948quantum,tomonaga1946relativistically, feynman2018space, feynman2018theory,feynman1950mathematical, dyson1949radiation, dyson1949s}、强相互作用\cite{yang1954conservation}及电弱相互作用\cite{glashow1961partial, salam1964electromagnetic, weinberg1967model}的不断深入研究，人类对自然界基本相互作用的认识持续深化。在几代卓越物理学家的努力下建立的粒子物理标准模型，不仅成功解释了已有实验现象，还预言了新夸克\ct{aubert1974experimental, abrams1974discovery, herb1977observation, innes1977observation, abe1995observation}、中间矢量玻色子\cite{barber1979discovery, arnison1983experimental, arnison1983experimental01}、第三代带电轻子及中微子\cite{perl1975evidence,kodama2001observation}和希格斯粒子\cite{aad2012observation}等新粒子，并通过后续实验不断得到验证。

粒子物理标准模型取得了巨大的成功，但依然对一些现象无法给出满意的解释。比如，为什么粒子的质量差别如此之大、为什么存在三代夸克和轻子、 为什么存在CP破坏、 为什么存在物质和反物质不对称、 为什么存在暗物质暗能量、 为什么存在引力且如此之弱等等。

\tb{追求统一，简洁和对自然界的更深刻的认识，是人类对于自然界的探究亘古不变的主旋律}。标准模型美而残缺，它的残缺之处，正是人类对于自然界的探究的动力。寻找新物理，主要包含四个方向：高亮度前沿（寻找新现象），高能量前沿（寻找新粒子），宇宙学前沿（寻找天外来客）和理论前沿（寻找新架构）。

重味物理是高亮度前沿研究的重要平台，专注于底夸克和粲夸克及其相互作用的研究。过去几十年，实验上关于重味强子弱衰变积累了大量的数据，同时重味物理具有丰富的物理可观测量，这使得精确研究成为可能。它能帮助我们更好地认识强相互作用的本质，精确检验粒子物理标准模型从而寻找寻找超出标准模型的新物理。

相比于B介子和底重子弱衰变，粲味强子弱衰变有其独特价值。作为唯一一类由上行夸克引发的弱衰变，这类过程对可能存在的高能标处的新物理高度敏感；同时该过程的典型能标接近非微扰区域，理解这类过程的动力学机制对我们理解强相互作用非微扰动力学机制具有启示作用。

区分所有末态粒子的过程被称为遍举过程\ct{Qin:2013tje}，而区分部分末态粒子的过程被称为单举过程\ct{Bigi:1993fe}。相比于遍举过程，单举过程基于算符乘积展开，可以从第一性原理出发系统地求解幂次修正的形式，这使得理论误差相对可控\ct{Neubert:1993mb}。与此同时，基于重夸克有效理论的算符乘积展开在底强子弱衰变领域取得了巨大的成功。然而，关于粲味强子弱衰变的理论研究相对较少，且幂次修正和微扰修正的收敛性较差，使得该类过程的研究相比底强子弱衰变更具挑战性。因此，深入研究此类过程对于该领域理论的进一步发展具有重要意义

目前，在重味物理领域仍存在一些长期未解的难题，其中一些包括：
\begin{itemize}
\item 粲强子寿命疑难：粲强子寿命的理论预言和实验测量结果存在较大差异，甚至$\mrm{D}^{+}$介子寿命在某些参数方案下的中心值小于零\ct{King:2021xqp, Cheng:2021qpd, Gratrex:2022xpm}。此外，粲重子寿命的理论预言对幂次修正和微扰修正非常敏感\ct{Cheng:2021qpd, Gratrex:2022xpm, Cheng:2023jpz}。
\item $\mrm{V_{cb}}$和$\mrm{V_{cb}}$疑难: 从单举衰变中抽取的CKM矩阵元$\mrm{V_{cb}}$和$\mrm{V_{cb}}$和从遍举衰变中的抽取结果相差3倍$\sigma$\ct{Gambino:2020jvv}。与其他类似的CKM矩阵元$\mrm{V_{cs}}$和$\mrm{V_{cd}}$进行对比是必要的，若从单举衰变中抽取的$\mrm{V_{cs}}$和$\mrm{V_{cd}}$与从遍举衰变中抽取的$\mrm{V_{cs}}$和$\mrm{V_{cd}}$一致，那么底强子单举弱衰变或许存在未知的新物理效应。
\item B介子反常: $\mrm{b\to s l^{+}l^{-}}$过程中发现测量得到的角分布和粒子物理标准模型的预言不符\ct{Archilli:2017xmu}，除了继续对$\mrm{b\to s l^{+}l^{-}}$进行高精度计算.更重要的是和其他类似的过程进行检验，特别是唯一上行夸克弱衰变$\mrm{c\to l^{+}l^{-}}$.
\end{itemize}

\tb{对于更清楚的认识，理解甚至解决这些疑难，粲味强子弱衰变恰好起到关键性的作用。}

由于粲味强子弱衰变的重要性，美国康奈尔大学的CLEO(Cornell Electron-Position storage Ring Large Dector)合作组和BESIII(Beijing Sepctrometer III)合作组对其进行了多次测量。2009年，CLEO合作组首次测量了$\mrm{D}^{0}$，$\mrm{D}^{+}$和$\mrm{D}_{\mrm{s}}^+$介子的半轻单举衰变分支比和电子能谱\ct{CLEO:2009uah}, 我国的BESIII合作组对$\mrm{D}^+_s$介子的半轻单举衰变分支比和电子能谱的测量结果进行更新\ct{BESIII:2021duu}。至此三种D介子的衰变分支比均具有较高的精度。对于粲重子$\Lambda_{c}$，BESIII合作组在2018年首次测量了其半轻单举衰变分支比和电子能谱\ct{BESIII:2018mug}，随后于2023年对结果进一步更新\ct{BESIII:2022cmg}，这为抽取非微扰强子矩阵元提供了必要的数据。未来基于神经网络技术，实验上能够实验X末态的完全重建以区分$\mrm{X}_{\mrm{d}}$和$\mrm{X}_{\mrm{s}}$末态。这将为从单举衰变中抽取的CKM矩阵元$\mrm{V}_{\mrm{cs}}$和$\mrm{V}_{\mrm{cd}}$提供必要的数据。同时，$\mrm{q}^{2}$谱也可以被提取。由于重参数不变性约束\ct{Luke:1992cs}，$\mrm{q}^{2}$矩依赖于较少的非微扰强子矩阵元，从而能够更高效地提取这些矩阵元，并与从电子能量矩中提取的非微扰矩阵元进行对比验证\ct{Mannel:2018mqv}。


粲强子寿命的理论计算依赖于一些非微扰的强子矩阵元, 其中一部分基于重夸克对称性从底强子的非微扰矩阵元演化而来，另一部分基于D和$\mrm{D}^{*}$质量的关系得到，它们的精度很低\ct{Neubert:1993mb, Uraltsev:2001ih, Falk:1992wt}。此外由于幂次修正的收敛性较差，目前的工作主要集中于树图水平上高量纲算符的计算\cite{Pak:2008qt, deBoer:2017way}。

综上所述，实验上已经对粲味强子单举衰变进行测量，并且对于某些衰变道的测量精度达到很高的水平，未来将对更多衰变道进行更精确的测量。，理论上对粲味强子单举衰变的研究仍然相对有限，尤其是系统的唯象学分析。本文首次对D介子半轻单举衰变做系统性的理论和唯象学研究，抽取并计算各类物理可观测量，非微扰参数的拟合估计精度与底强子的非微扰参数相一致，并讨论了质量方案对微扰修正和幂次修正的收敛性影响，提出1S质量方案是目前计算粲强子寿命过程中最佳的质量方案。该工作对粲味强子弱衰变的理论研究具有重要意义，同时也有助于理解甚至解决粲强子寿命疑难。此外, 对从单举衰变中抽取CKM矩阵元$\mrm{V}_{\mrm{cs}}$,$\mrm{V}_{\mrm{cd}}$以及粲重子非微扰参数的拟合估计奠定了理论基础。

本文的主要分为以下几个部分，第二章是对粒子物理标准模型的介绍，主要包括和本文研究内容密切相关的量子色动力学，重夸克有效理论和基于重夸克有效理论的局域算符乘积展开。第三章基于重夸克有效理论和算符乘积展开对粲介子半轻单举衰变进行系统研究，包括精确至两圈图水平的微扰修正，精确至量纲六算符的幂次修正。第四章基于已有的实验数据，抽取相应的物理观测量进行全局拟合。第五章是总结与展望，对本文的工作进行总结并指出不足之处，同时对未来做短期和长期的展望。